Ist die Sprache
\[
L
=
\{
\texttt{a}^n
\texttt{b}^m
\texttt{c}^m
\texttt{d}^n
\mid
n>0\wedge m>0
\}
\]
über dem Alphabet
$\Sigma = \{
\texttt{a},
\texttt{b},
\texttt{c},
\texttt{d}\}$
kontextfrei?

\begin{loesung}
Die Sprache ist kontextfrei, denn Sie kann mit der Grammatik
\begin{align*}
S&\to \texttt{a}S\texttt{d} \mid T \\
T&\to \texttt{b}T\texttt{c} \mid \texttt{bc}
\end{align*}
erzeugt wird.

Alternativ kann man sie auch mit dem Stackautomaten
\begin{center}
\def\l{2.8}
\begin{tikzpicture}[>=latex,thick]
\coordinate (s) at ({-\l+0.5},0);
\coordinate (q0) at (0,0);
\coordinate (q1) at ({1*\l},0);
\coordinate (q2) at ({2*\l},0);
\coordinate (q3) at ({3*\l},0);
\coordinate (q4) at ({4*\l},0);
\coordinate (e) at ({5*\l},0);
\node at (q0) {$q_0$};
\node at (q1) {$q_1$};
\node at (q2) {$q_2$};
\node at (q3) {$q_3$};
\node at (q4) {$q_4$};
\node at (e) {$e$};
\draw (q0) circle[radius=0.35];
\draw (q1) circle[radius=0.35];
\draw (q2) circle[radius=0.35];
\draw (q3) circle[radius=0.35];
\draw (q4) circle[radius=0.35];
\draw (e) circle[radius=0.35];
\draw (e) circle[radius=0.30];
\draw[->,shorten >= 0.35cm] (s) -- (q0);
\draw[->,shorten <= 0.35cm,shorten >= 0.35cm] (q0) -- (q1);
\node at ($0.5*(q0)+0.5*(q1)$) [above]
	{$\varepsilon,\varepsilon\to\texttt{\$}\mathstrut$};
\draw[->,shorten <= 0.35cm,shorten >= 0.35cm] (q1) -- (q2);
\node at ($0.5*(q1)+0.5*(q2)$) [above]
	{$\texttt{a},\varepsilon\to\texttt{a}\mathstrut$};
\draw[->,shorten <= 0.35cm,shorten >= 0.35cm] (q2) -- (q3);
\node at ($0.5*(q2)+0.5*(q3)$) [above]
	{$\texttt{b},\varepsilon\to\texttt{b}\mathstrut$};
\draw[->,shorten <= 0.35cm,shorten >= 0.35cm] (q3) -- (q4);
\node at ($0.5*(q3)+0.5*(q4)$) [above]
	{$\varepsilon,\varepsilon\to\varepsilon\mathstrut$};
\draw[->,shorten <= 0.35cm,shorten >= 0.35cm] (q4) -- (e);
\node at ($0.5*(q4)+0.5*(e)$) [above]
	{$\varepsilon,\texttt{\$}\to\varepsilon\mathstrut$};
\draw[->,shorten <= 0.35cm,shorten >= 0.35cm]
	(q1) to[out=60,in=120,distance=1.2cm] (q1);
\node at ($(q1)+(0,0.8)$) [above] {$\texttt{a},\varepsilon\to\texttt{a}\mathstrut$};
\draw[->,shorten <= 0.35cm,shorten >= 0.35cm]
	(q2) to[out=60,in=120,distance=1.2cm] (q2);
\node at ($(q2)+(0,0.8)$) [above] {$\texttt{b},\varepsilon\to\texttt{b}\mathstrut$};
\draw[->,shorten <= 0.35cm,shorten >= 0.35cm]
	(q3) to[out=60,in=120,distance=1.2cm] (q3);
\node at ($(q3)+(0,0.8)$) [above] {$\texttt{c},\texttt{b}\to\varepsilon\mathstrut$};
\draw[->,shorten <= 0.35cm,shorten >= 0.35cm]
	(q4) to[out=60,in=120,distance=1.2cm] (q4);
\node at ($(q4)+(0,0.8)$) [above] {$\texttt{d},\texttt{a}\to\varepsilon\mathstrut$};
\end{tikzpicture}
\end{center}
akzeptieren.
Die Übergänge von $q_1$ nach $q_2$ und von $q_2$ nach $q_3$ stellen sicher,
dass mindestens ein \texttt{a} bzw.~mindestens ein \texttt{b} auf den
Stack gelegt wird.
\end{loesung}

\begin{bewertung}
Teilgrammatik für die Teile aus \texttt{a} und \texttt{d} ({\bf A}) 1 Punkt,
Schachtelungstrategie ({\bf S}) 1 Punkt,
Verkettung der Grammatiken ({\bf V}) 1 Punkt,
Teilgrammatik für die Teile aus \texttt{b} und \texttt{c} ({\bf B}) 1 Punkt,
Terminierung mit \texttt{bc} ({\bf T}) 1 Punkt,
Schlussfolgerung kontextfrei ({\bf K}) 1 Punkt.
\end{bewertung}

